\documentclass[12pt,a4paper]{article}

\usepackage[utf8]{inputenc}
\usepackage[T1]{fontenc}
\usepackage[turkish,english]{babel}
\usepackage{geometry}
\usepackage{amsmath,amssymb}
\usepackage{booktabs}
\usepackage{longtable}
\usepackage{array}
\usepackage{enumitem}
\usepackage{hyperref}
\usepackage{tikz}
\usepackage{pgfplots}
\usepackage{xcolor}
\pgfplotsset{compat=1.18}
\usetikzlibrary{shapes.geometric,arrows.meta,positioning,calc,patterns}

% Renk paleti
\definecolor{cRemote}{HTML}{4E79A7}
\definecolor{cContext}{HTML}{F28E2B}
\definecolor{cInsitu}{HTML}{E15759}
\definecolor{cChem}{HTML}{76B7B2}
\definecolor{cContam}{HTML}{59A14F}
\definecolor{cFusion}{HTML}{EDC948}
\definecolor{cLow}{HTML}{D62728}
\definecolor{cMid}{HTML}{FF7F0E}
\definecolor{cHigh}{HTML}{2CA02C}
\definecolor{cExtra}{HTML}{9467BD}

\geometry{margin=2.5cm}
\setlength{\parskip}{0.6em}
\setlength{\parindent}{0pt}
\raggedbottom

\title{Derin Uzay Astrobiyolojisi İçin Biyo-iz Güven Skoru:\\
Uzaktan Algılamadan Yakın Doğrulamaya Çok Modlu Karar Modeli}
\author{Sercan Demirhan \and Gizem Demirhan}
\date{\today}

\begin{document}
\selectlanguage{turkish}
\shorthandoff{=}
\maketitle

\begin{abstract}
Bu çalışma, Ay ve Mars odaklı astrobiyoloji görevlerinde yaşam izlerinin güvenilir biçimde değerlendirilmesi için çok modlu bir karar modeli önermektedir. Yaklaşım, uzaktan algılama (teleskop/spektral), jeolojik bağlam analizi, yüzey morfolojisi, kimyasal analiz ve izotopik doğrulama adımlarını tek bir biyo-iz güven skorunda birleştirir. Model, Bayes güncelleme, makine öğrenmesi tabanlı sınıflandırma ve kontaminasyon riski cezalandırma mekanizmasını bir arada kullanır. Çıktı olarak sahaya özgü karar etiketleri, operasyonel örnekleme öncelikleri ve 90 günlük uygulanabilir asgari ürün (MVP) yol haritası sunulur.
\end{abstract}

\section{Giriş}
Astrobiyolojide en kritik sorunlardan biri, biyotik ve abiyotik süreçlerin benzer gözlemsel izler üretebilmesidir. Bu nedenle tek bir ölçüm kanalına dayalı yorumlar yüksek yanlış pozitif riski taşır. Bu makale, ``çoklu kanıt'' ilkesine dayanan bir karar çerçevesi önermektedir \cite{allwood2006}: morfoloji, mineralojik-jeolojik bağlam, organik kimya, izotopik fraksiyonasyon \cite{schidlowski2001} ve kalite-kontaminasyon göstergeleri birlikte değerlendirilir.

\section{Araştırma Amacı ve Hipotez}
\subsection{Amaç}
Ay ve Mars görevlerinde yaşam benzeri süreçlere ait olası biyo-izleri, uzak ölçümden yakın doğrulamaya uzanan bir zincirde nicel olarak puanlamak.

\subsection{Ana Hipotez}
Morfolojik benzerlik tek başına yeterli değildir. Güvenilir bir yorum için şu üç eksen birlikte tutarlı olmalıdır:
\begin{enumerate}[label=\arabic*.]
    \item Jeolojik ve mineralojik bağlam,
    \item Kimyasal ve izotopik imzalar,
    \item Kontaminasyon ve abiyotik taklit risklerinin dışlanması.
\end{enumerate}

\section{Karar Modeli ve Matematiksel Çerçeve}
\subsection{Temel Olasılık Modeli}
Nihai hedef skor aşağıdaki sonsal (\textit{posterior}) olasılık olarak tanımlanır:
\begin{equation}
S_{bio} = P(\text{ya\u015fam benzeri s\u00fcre\u00e7} \mid X_{orbital}, X_{surface}, X_{chem}, X_{iso}, X_{qc}).
\end{equation}

\subsubsection*{Denklem Açıklaması}
Burada $X_{orbital}$ uzaktan algılama spektral/jeomorfolojik gözlemlerini, $X_{surface}$ yüzey morfoloji ölçümlerini, $X_{chem}$ kimyasal analiz çıktısını, $X_{iso}$ izotopik kanıtları ve $X_{qc}$ kalite-kontaminasyon denetimlerini temsil eder. Dolayısıyla $S_{bio}\in[0,1]$ değeri, tüm kanıtlar birlikte değerlendirildiğinde ``biyotik sürece ait olma'' inancının güncel olasılık değeridir.

Modelin Bayes açılımı iki rakip hipotez üzerinden yazılabilir:
\begin{equation}
P(H_B\mid D)=\frac{P(D\mid H_B)P(H_B)}{P(D\mid H_B)P(H_B)+P(D\mid H_A)P(H_A)},
\end{equation}
burada $H_B$ biyotik hipotezi, $H_A$ abiyotik hipotezi ve $D=\{X_{orbital},X_{surface},X_{chem},X_{iso},X_{qc}\}$ birleşik gözlem vektörüdür. Bu formülasyonda $P(H_B)$ öncül (\textit{prior}) inancı, $P(D\mid H_B)$ biyotik durumda gözlemlerin olabilirliğini, payda terimi ise normalizasyon sabiti olarak toplam kanıt olasılığını ifade eder.

Pratik yorum: $P(D\mid H_A)$ terimi büyüdükçe abiyotik açıklama gücü artar ve $S_{bio}$ düşer; bu sayede model yalnızca ``biyotik kanıt toplamakla'' kalmaz, aynı zamanda yanlış pozitif üreten abiyotik süreçleri de nicel olarak baskılar.

\subsection{Uygulama İçin Birleşik Skor}
Operasyonel kullanımda ağırlıklı skor formu:
\begin{equation}
S_{final} = \alpha S_{remote} + \beta S_{context} + \gamma S_{in\text{-}situ} + \delta S_{chem\text{-}iso} - \lambda R_{contam}.
\end{equation}

Burada:
\begin{itemize}
    \item $S_{remote}$: Uzaktan algılama tarama skoru,
    \item $S_{context}$: Jeolojik bağlam uygunluk skoru,
    \item $S_{in\text{-}situ}$: Yüzey morfoloji skoru,
    \item $S_{chem\text{-}iso}$: Kimyasal ve izotopik doğrulama skoru,
    \item $R_{contam}$: Kontaminasyon ve abiyotik taklit risk terimi ($R_{contam}\in[0,1]$; temiz numunede $\approx 0$).
\end{itemize}
Ağırlıklar $\alpha+\beta+\gamma+\delta=1{,}00$ olacak şekilde normalize edilmiştir. $\lambda$ çıkarma katsayısı olup pratikte $R_{contam}$ değeri çoğu senaryoda düşüktür; bu nedenle $S_{final}\in[0,1]$ aralığında kalır.

\subsection{Aşamalı Bayes Güncelleme}
Her yeni ölçüm kümesiyle sonsal olasılık güncellenir:
\begin{equation}
P(L \mid D_{1:k}) = \frac{P(D_k \mid L, D_{1:k-1})\,P(L \mid D_{1:k-1})}{P(D_k \mid D_{1:k-1})},
\end{equation}
burada $L$ yaşam benzeri süreci, $D_{1:k}$ ise sırasıyla gelen veri paketlerini ifade eder.

\subsubsection*{Denklem Bileşenlerinin Yorumu}
\begin{itemize}
    \item $P(L \mid D_{1:k-1})$: bir önceki adımın sonsal değeri; yeni adım için öncül (\textit{prior}) olarak kullanılır.
    \item $P(D_k \mid L, D_{1:k-1})$: $k$. veri paketinin, yaşam hipotezi doğruyken gözlenme olabilirliği (\textit{likelihood}).
    \item $P(D_k \mid D_{1:k-1})$: normalizasyon terimi (kanıt olasılığı); sonucun $[0,1]$ aralığında kalmasını sağlar.
    \item $P(L \mid D_{1:k})$: yeni veriyi dahil ettikten sonraki güncel güven düzeyi (sonsal olasılık).
\end{itemize}

Bu güncelleme kuralı, işlem hattındaki her fazı (uzaktan tarama, bağlam, morfoloji, kimya-izotop) ardışık kanıt olarak ele alır. Eğer yeni veri paketi $L$ ile yüksek uyumluysa ($P(D_k\mid L,\cdot)$ yüksek), skor artar; veri abiyotik açıklamalarla daha tutarlıysa skor azalır. Böylece model tek bir ölçüme değil, zaman içinde biriken çoklu kanıta dayalı karar verir.

Operasyonel olarak her adım sonunda hesaplanan $P(L\mid D_{1:k})$ değeri, bir sonraki ölçüm ihtiyacını belirlemek için eşiklerle karşılaştırılır: düşük değerlerde ek ölçüm planlanır, yüksek ve kararlı değerlerde hedef önceliği artırılır.

\section{Yöntem: Uzaktan Yakına İşlem Hattı}
\subsection{Faz 0: Etiket ve Referans Kütüphanesi}
\begin{itemize}
    \item Dünya analogları: hidrotermal, evaporitik, biyofilmli ve abiyotik sahalar \cite{kokaly2017}.
    \item Sınıflar: biyotik benzeri, abiyotik benzeri, belirsiz.
    \item Çıktı: Eğitim verisi şeması ve kalite protokolü.
\end{itemize}

\subsection{Faz 1: Uzaktan Tarama (Teleskop/Spektroskopi)}
\begin{itemize}
    \item Girdi: VNIR-SWIR, termal bantlar, topografya, morfometri \cite{clark1984}.
    \item Model: 1D-CNN veya spektral \textit{transformer}.
    \item Ek model: Jeomorfolojik segmentasyon için CNN tabanlı nesne tanıma.
    \item Çıktı: Adaylık haritası ve $S_{remote}$.
\end{itemize}

\subsection{Faz 2: Jeolojik Bağlam Skoru}
\begin{itemize}
    \item Girdi: Delta/yelpaze morfolojisi, paleolaküstrin birimler, kil-sülfat-silika varlığı.
    \item Model: \textit{Gradient boosting} veya Bayes ağları.
    \item Çıktı: $S_{context}$.
\end{itemize}

\subsection{Faz 3: Yüzey Morfoloji Analizi}
\begin{itemize}
    \item Girdi: Makro/mikro görüntü, 3B topoloji, doku metrikleri.
    \item Hedef: Stromatolit benzeri kubbe/sütun, filamenter doku, ritmik laminasyon.
    \item Model: Çok ölçekli CNN + öznitelik tabanlı istatistik.
    \item Çıktı: $S_{in\text{-}situ}$.
\end{itemize}

\subsection{Faz 4: Kimyasal ve İzotopik Doğrulama}
\begin{itemize}
    \item Girdi: GC-MS, LC-MS, Raman, FTIR, IRMS, redoks sensörleri.
    \item Model: Çok modlu geç birleştirme (\textit{late fusion}) + kalibre edilmiş lojistik katman.
    \item Çıktı: $S_{chem\text{-}iso}$.
\end{itemize}

\subsection{Faz 5--6: Karar Füzyonu ve Operasyonel Planlama}
\begin{itemize}
    \item Bayes güncelleme ile nihai skor.
    \item Belirsizlik ve veri eksikliği cezası.
    \item Karar etiketleri: düşük, orta, yüksek güven, ek ölçüm gerekli.
\end{itemize}

% ===== ŞEKİL 1: İşlem Hattı Akış Diyagramı =====
\begin{figure}[ht]
\centering
\begin{tikzpicture}[
    node distance=0.8cm and 0.4cm,
    block/.style={rectangle, draw, rounded corners, minimum width=2.8cm, minimum height=0.9cm, font=\footnotesize\bfseries, text=white, align=center},
    arrow/.style={-{Stealth[length=3mm]}, thick},
    scorelbl/.style={font=\tiny, text=black, anchor=west}
]
% Faz kutuları
\node[block, fill=cRemote]  (f0) {Faz 0\\Referans Kütüph.};
\node[block, fill=cRemote, right=of f0]  (f1) {Faz 1\\Uzaktan Tarama};
\node[block, fill=cContext, right=of f1] (f2) {Faz 2\\Jeolojik Bağlam};
\node[block, fill=cInsitu, right=of f2]  (f3) {Faz 3\\Morfoloji};
\node[block, fill=cChem, right=of f3]    (f4) {Faz 4\\Kimya-İzotop};
\node[block, fill=cContam, below=1.2cm of f2, xshift=1.6cm, minimum width=11.22cm] (f56) {Faz 5--6: Karar Füzyonu};
\node[block, fill=cFusion, below=0.8cm of f56, minimum width=3.5cm, text=black] (out) {$S_{final}$ + Etiket};

% Oklar
\draw[arrow] (f0) -- (f1);
\draw[arrow] (f1) -- (f2);
\draw[arrow] (f2) -- (f3);
\draw[arrow] (f3) -- (f4);

% Dikey oklar füzyon'a
\draw[arrow, cRemote]  (f1.south) -- ++(0,-0.35) -| (f56.north -| f1)  node[scorelbl, pos=0.25, left]{$S_{remote}$};
\draw[arrow, cContext] (f2.south) -- ++(0,-0.35) -| (f56.north -| f2) node[scorelbl, pos=0.25, left]{$S_{context}$};
\draw[arrow, cInsitu]  (f3.south) -- ++(0,-0.35) -| (f56.north -| f3) node[scorelbl, pos=0.25, right]{$S_{in\text{-}situ}$};
\draw[arrow, cChem]    (f4.south) -- ++(0,-0.35) -| (f56.north -| f4) node[scorelbl, pos=0.25, right]{$S_{chem}$};

\draw[arrow] (f56) -- (out);
\end{tikzpicture}
\caption{Biyo-iz güven skoru işlem hattı: uzaktan taramadan karar füzyonuna.}
\label{fig:pipeline}
\end{figure}

\section{Astrobiyolojik Analiz Kategorileri}

% ===== ŞEKİL 7: Dünya-Mars Analog Karşılaştırması =====
\begin{figure}[ht]
\centering
\begin{tikzpicture}[
    panel/.style={draw=gray!60, thick, rounded corners=2pt, minimum width=6.8cm, minimum height=6.2cm, anchor=north west},
    plabel/.style={font=\scriptsize\bfseries, fill=white, inner sep=2pt},
    annot/.style={font=\tiny, text=gray!70!black},
    sublabel/.style={font=\tiny\bfseries, fill=white, inner sep=1.5pt, rounded corners=1pt, anchor=west},
]

% ---- Sol panel: Dünya Analogları ----
\node[panel, fill=cRemote!5] (earthP) at (-1.0,0) {};
\node[plabel, anchor=north west] at (-0.9,-0.1) {Dünya Analogları};

% Shark Bay stromatolit şeması
\begin{scope}[shift={(0.4,-2.0)}]
  \node[sublabel] at (0,0.85) {(a) Shark Bay, Avustralya};
  % Deniz seviyesi
  \draw[cRemote, thick] (0,0) -- (5,0);
  \node[annot, anchor=east] at (0,0.1) {\textit{deniz sev.}};
  % Stromatolit kubbeleri
  \foreach \x/\h in {0.6/0.45, 1.3/0.55, 2.0/0.40, 2.7/0.50, 3.4/0.35, 4.1/0.48} {
    \fill[cHigh!40!cContext, rounded corners=3pt] (\x-0.25, -0.15) rectangle (\x+0.25, \h);
    \draw[cHigh!70!black, thick, rounded corners=3pt] (\x-0.25, -0.15) rectangle (\x+0.25, \h);
  }
  % Laminasyon çizgileri (stromatolit iç yapısı)
  \foreach \x/\h in {0.6/0.45, 1.3/0.55, 2.0/0.40, 2.7/0.50, 3.4/0.35, 4.1/0.48} {
    \foreach \dy in {0.08, 0.18, 0.28} {
      \draw[cHigh!50!black, very thin] (\x-0.18, \dy) -- (\x+0.18, \dy);
    }
  }
  % Alt sediman
  \fill[brown!20] (0,-0.15) rectangle (5,-0.55);
  \draw[brown!40, densely dashed] (0,-0.35) -- (5,-0.35);
  \node[annot] at (2.5,-0.35) {karbonat sediman};
  % Ölçek
  \draw[|-|, thick] (4.4, -0.80) -- (5, -0.80) node[midway, below, font=\tiny] {$\sim$0{,}5\,m};
\end{scope}

% Yellowstone sinter şeması
\begin{scope}[shift={(0.4,-4.2)}]
  \node[sublabel] at (0,0.65) {(b) Yellowstone, ABD};
  % Sıcak kaynak havuzu
  \fill[cRemote!25] (0.2,0) .. controls (1,-0.25) and (2,-0.25) .. (2.5,0)
                     -- (2.5,-0.3) -- (0.2,-0.3) -- cycle;
  \draw[cRemote!60, thick] (0.2,0) .. controls (1,-0.25) and (2,-0.25) .. (2.5,0);
  \node[annot] at (1.35,-0.15) {\textit{termal havuz}};
  % Sinter terası
  \fill[cContext!30] (2.5,0) -- (3.0,0.15) -- (3.5,0.25) -- (4.2,0.30) -- (5,0.20)
                     -- (5,-0.3) -- (2.5,-0.3) -- cycle;
  \draw[cContext!70, thick] (2.5,0) -- (3.0,0.15) -- (3.5,0.25) -- (4.2,0.30) -- (5,0.20);
  % Mikrobiyal mat çizgileri
  \foreach \y in {0.0, -0.10, -0.20} {
    \draw[cHigh!60, thin] (2.6,\y) -- (4.8, {\y+0.22});
  }
  \node[annot] at (4.0,0.0) {silika sinter};
  % Buhar
  \draw[gray!40, thick, decorate, decoration={snake, amplitude=1.5pt, segment length=5pt}]
    (1.0, 0.05) -- (1.0, 0.45);
  \draw[gray!40, thick, decorate, decoration={snake, amplitude=1.5pt, segment length=5pt}]
    (1.8, 0.0) -- (1.8, 0.40);
\end{scope}

% ---- Sağ panel: Mars Hedef Bölgeleri ----
\node[panel, fill=cLow!5] (marsP) at (6.9,0) {};
\node[plabel, anchor=north west] at (7.0,-0.1) {Mars Hedef Bölgeleri};

% Jezero delta şeması
\begin{scope}[shift={(7.8,-2.0)}]
  \node[sublabel] at (0,0.65) {(c) Jezero Krateri Deltas{\i}};
  % Krater kenarı
  \draw[cLow!70, thick] (0,0.15) -- (0.3,-0.1) -- (0.5,-0.15);
  % Delta yelpaze
  \fill[cContext!35] (0.5,-0.15) -- (1.2,0.20) -- (2.0,0.30) -- (3.0,0.25) -- (3.8,0.10)
                     -- (3.8,-0.45) -- (0.5,-0.45) -- cycle;
  \draw[cContext!70, thick] (0.5,-0.15) -- (1.2,0.20) -- (2.0,0.30) -- (3.0,0.25) -- (3.8,0.10);
  % Delta iç katmanlar
  \draw[cContext!50, thin, dashed] (0.7,-0.10) -- (3.6, 0.05);
  \draw[cContext!50, thin, dashed] (0.6,-0.25) -- (3.7, -0.10);
  \draw[cContext!50, thin, dashed] (0.6,-0.35) -- (3.7, -0.25);
  % Mineral notları
  \node[annot, fill=cChem!15, rounded corners=1pt, inner sep=1pt] at (1.5, 0.0) {smektit};
  \node[annot, fill=cContext!15, rounded corners=1pt, inner sep=1pt] at (2.8, -0.12) {karbonat};
  % Krater tabanı
  \draw[cLow!50, thick] (3.8,0.10) -- (4.3,-0.10) -- (5.0,-0.15);
  \fill[cLow!10] (0.5,-0.45) rectangle (5,-0.55);
  \node[annot] at (2.5,-0.50) {bazalt taban};
  % Kanal giriş oku
  \draw[-{Stealth}, cRemote, thick] (-0.1, 0.25) -- (0.5, 0.0) node[midway, above, annot]{kanal};
\end{scope}

% Mars sülfat damar şeması
\begin{scope}[shift={(7.8,-4.2)}]
  \node[sublabel] at (0,0.65) {(d) Mars Sülfat Damarları};
  % Kayaç matris
  \fill[cLow!12] (0,-0.05) rectangle (5,-0.65);
  \draw[cLow!30] (0,-0.05) -- (5,-0.05);
  % Bazalt dokusu (rastgele çizgiler)
  \foreach \x in {0.3, 1.0, 1.7, 2.5, 3.2, 3.9, 4.5} {
    \draw[cLow!20, very thin] (\x,-0.15) -- ({\x+0.3},-0.55);
  }
  % Sülfat damarları (parlak beyaz çizgiler)
  \draw[white, line width=2.5pt] (0.4,-0.10) -- (0.8,-0.60);
  \draw[cFusion!80, line width=1.5pt] (0.4,-0.10) -- (0.8,-0.60);
  \draw[white, line width=2.5pt] (2.0,-0.08) -- (2.5,-0.62);
  \draw[cFusion!80, line width=1.5pt] (2.0,-0.08) -- (2.5,-0.62);
  \draw[white, line width=2.5pt] (3.5,-0.12) -- (3.2,-0.58);
  \draw[cFusion!80, line width=1.5pt] (3.5,-0.12) -- (3.2,-0.58);
  \draw[white, line width=1.8pt] (1.2,-0.20) -- (1.8,-0.22) -- (2.4,-0.35);
  \draw[cFusion!70, line width=1pt] (1.2,-0.20) -- (1.8,-0.22) -- (2.4,-0.35);
  % Etiketler
  \node[annot, fill=cFusion!20, rounded corners=1pt, inner sep=1pt] at (1.0,-0.35) {CaSO$_4$};
  \node[annot] at (4.2,-0.35) {bazalt matris};
  % Ölçek
  \draw[|-|, thick] (4.0, -0.90) -- (5, -0.90) node[midway, below, font=\tiny] {$\sim$5\,cm};
\end{scope}

% Ortadaki karşılaştırma okları
\draw[{Stealth}-{Stealth}, thick, cExtra, dashed] (6.4,-2.2) -- (6.8,-2.2)
  node[midway, above, font=\tiny, text=cExtra, rotate=90, yshift=3pt] {analog};
\draw[{Stealth}-{Stealth}, thick, cExtra, dashed] (6.4,-4.5) -- (6.8,-4.5)
  node[midway, above, font=\tiny, text=cExtra, rotate=90, yshift=3pt] {analog};

\end{tikzpicture}
\caption{Dünya analog sahaları ile Mars hedef bölgelerinin şematik karşılaştırması. (a)~Shark Bay stromatolitleri: karbonat sediman üzerinde ritmik laminasyonlu kubbe yapılar \cite{allwood2006}. (b)~Yellowstone sıcak kaynakları: silika sinter terasları ve mikrobiyal mat katmanları. (c)~Jezero krateri deltası: smektit kil ve karbonat minerallerince zengin paleolaküstrin çökel dizisi \cite{farley2020}. (d)~Mars sülfat (CaSO$_4$) damarları: bazalt matris içinde sulu alterasyon kanıtı.}
\label{fig:analog}
\end{figure}

\subsection{Jeolojik Oluşumların Analizi}

\subsubsection*{Ay}
Ay'da volkanik bazaltlar (\textit{mare}) ve anortozit kabuk, düşük su aktivitesi nedeniyle biyolojik potansiyeli sınırlar \cite{smith2010, robinson2010}. Regolit stratigrafisi; mikrometeorit bombardımanı ve cam kürecikler (aglütina) içermektedir. Kalıcı gölge bölgeler (PSR) ise krater içinde buz deposu oluşturabilmektedir.

\textbf{Yöntemler:}
\begin{itemize}
    \item Yüksek çözünürlüklü görüntüleme + fotogrametri: stratigrafi ve katman süreklilik analizi,
    \item Spektral haritalama (VNIR, SWIR): Fe$^{2+}$/Fe$^{3+}$ oranları, piroksen/olivin ayrımı,
    \item Nötron spektrometresi: hidrojen \textit{proxy}'si üzerinden su/buz tahmini,
    \item Yer radarı (GPR): gömülü buz ve lav tüpü tespiti,
    \item İzotop jeokronolojisi (U-Pb, Ar-Ar): yüzey yenilenme tarihleri.
\end{itemize}

\subsubsection*{Mars}
Mars için öncelikli hedefler \cite{farley2020}: delta/yelpaze sistemleri, paleolaküstrin çökeller, kil mineralleri (smektit), sülfatlar, silika zenginleşmeleri, hidrotermal alterasyon zonları, damar yapıları ve sedimanter laminasyon (mikrobiyal mat analoğu olabilecek ince tabakalar).

\textbf{Yöntemler:}
\begin{itemize}
    \item Orbital + yerinde jeomorfoloji: delta morfometrisi, kanal ağları \cite{mcewen2007, murchie2007},
    \item Mineraloji (VNIR/IR, Raman, XRD): kil/sülfat/silika faz tayini \cite{bhartia2021, maurice2021},
    \item Mikro doku analizi ($\mu$-görüntüleme): laminasyon ve stromatolit benzeri yapılar,
    \item Jeokimyasal haritalama (LIBS, XRF): element dağılımı ve redoks gradyanları \cite{allwood2020}.
\end{itemize}

\subsection{Görsel İşaretlerin (Morfolojik Biyo-iz) Analizi}

\textbf{Hedef işaretler:}
\begin{itemize}
    \item Mikrometre--milimetre ölçekli stromatolitik kubbeler ve sütunlar,
    \item Filamenter/koloni benzeri yapılar (biyofilmler),
    \item Sediman içinde ritmik laminasyon (günlük/mevsimsel döngüler).
\end{itemize}

\textbf{Yöntemler:}
\begin{itemize}
    \item Makro/mikro görüntüleme: stereo kameralar ve mikroskobik görüntüleyici,
    \item 3B rekonstrüksiyon: fotogrametri/LiDAR ile yüzey topolojisi,
    \item Doku metrikleri: fraktal boyut ve yönlenme dağılımı (abiyotik-biyotik ayrımı),
    \item Çapraz kesit inceleme: mikro tomografi ($\mu$CT) ile iç yapı sürekliliğinin doğrulanması,
    \item Kontrol karşılaştırmaları: evaporitik kristaller, rüzg\^ar oyukları gibi abiyotik analoglarla istatistiksel ayrım.
\end{itemize}

\subsection{Kimyasal ve Biyokimyasal İşaretlerin Analizi}

\textbf{Hedef işaretler:}
\begin{itemize}
    \item Organik moleküller: aromatikler, alifatik zincirler, potansiyel biyobelirteçler (lipid türevleri),
    \item İzotopik fraksiyonasyon: $\delta^{13}$C, $\delta^{34}$S, $\delta$D sapmaları (biyotik süreçler genellikle hafif izotop zenginleşmesi üretir) \cite{schidlowski2001},
    \item Redoks çiftleri: Fe$^{2+}$/Fe$^{3+}$, S$^{2-}$/SO$_4^{2-}$ dengesizlikleri,
    \item Perklorat/oksidan ortam etkileri (özellikle Mars'ta organik bozunma riski).
\end{itemize}

\textbf{Yöntemler:}
\begin{itemize}
    \item GC-MS\@ / Piroliz-GC-MS\@: uçucu ve yarı uçucu organikler,
    \item LC-MS\@: termal bozunmaya hassas bileşikler,
    \item Raman spektroskopisi\@: organik bantlar ve mineral bağlamı birlikte,
    \item FTIR\@: fonksiyonel grup tayini,
    \item İzotop oran kütle spektrometrisi (IRMS)\@: $\delta$ ölçümleri,
    \item Elektrokimya sensörleri\@: redoks profilleri ve oksidan yük.
\end{itemize}

% ===== ŞEKİL 8: Biyotik vs Abiyotik Spektral İmza Karşılaştırması =====
\begin{figure}[ht]
\centering
\begin{tikzpicture}
\begin{axis}[
    width=13cm, height=7cm,
    xlabel={Dalga boyu ($\mu$m)},
    ylabel={Yansıtma (normalize)},
    xmin=0.4, xmax=2.6,
    ymin=-0.08, ymax=1.05,
    xtick={0.5, 0.8, 1.0, 1.4, 1.9, 2.2, 2.5},
    xticklabel style={font=\scriptsize},
    yticklabel style={font=\scriptsize},
    xlabel style={font=\small},
    ylabel style={font=\small},
    grid=major,
    grid style={gray!15},
    legend style={at={(0.02,0.98)}, anchor=north west, font=\scriptsize,
                  draw=gray!40, fill=white, fill opacity=0.9},
    legend cell align={left},
]

% Stromatolit (biyotik) — tipik karbonat+organik imza
\addplot[cHigh, thick, smooth] coordinates {
    (0.40, 0.12) (0.50, 0.18) (0.55, 0.30) (0.60, 0.42)
    (0.67, 0.55) (0.75, 0.68) (0.85, 0.76) (1.00, 0.82)
    (1.10, 0.80) (1.20, 0.78) (1.40, 0.72) (1.50, 0.65)
    (1.60, 0.62) (1.70, 0.58) (1.80, 0.50) (1.90, 0.38)
    (2.00, 0.44) (2.10, 0.52) (2.20, 0.55) (2.30, 0.48)
    (2.35, 0.40) (2.40, 0.42) (2.50, 0.38)
};

% Smektit kil (jeolojik bağlam) — 1.4 µm ve 1.9 µm'de derin absorpsiyon
\addplot[cContext, thick, smooth, dashed] coordinates {
    (0.40, 0.15) (0.50, 0.25) (0.60, 0.38) (0.70, 0.50)
    (0.80, 0.60) (0.90, 0.68) (1.00, 0.75) (1.10, 0.78)
    (1.20, 0.76) (1.30, 0.70) (1.40, 0.42) (1.50, 0.60)
    (1.60, 0.68) (1.70, 0.72) (1.80, 0.65) (1.90, 0.30)
    (2.00, 0.55) (2.10, 0.62) (2.20, 0.56) (2.30, 0.35)
    (2.35, 0.40) (2.40, 0.38) (2.50, 0.35)
};

% Demir oksit (abiyotik taklitçi) — kırmızı eğimli, 0.9 µm'de geniş absorpsiyon
\addplot[cLow, thick, smooth, densely dotted] coordinates {
    (0.40, 0.05) (0.45, 0.07) (0.50, 0.10) (0.55, 0.18)
    (0.60, 0.30) (0.65, 0.45) (0.70, 0.55) (0.75, 0.60)
    (0.80, 0.55) (0.85, 0.48) (0.90, 0.40) (0.95, 0.42)
    (1.00, 0.50) (1.10, 0.60) (1.20, 0.65) (1.40, 0.70)
    (1.60, 0.72) (1.80, 0.73) (1.90, 0.70) (2.00, 0.72)
    (2.10, 0.73) (2.20, 0.74) (2.30, 0.73) (2.50, 0.72)
};

% Evaporit (abiyotik — düz profil)
\addplot[cExtra, thick, smooth, dashdotted] coordinates {
    (0.40, 0.55) (0.50, 0.62) (0.60, 0.70) (0.70, 0.78)
    (0.80, 0.84) (0.90, 0.88) (1.00, 0.90) (1.10, 0.91)
    (1.20, 0.90) (1.40, 0.82) (1.50, 0.78) (1.60, 0.80)
    (1.70, 0.82) (1.80, 0.80) (1.90, 0.72) (2.00, 0.78)
    (2.10, 0.82) (2.20, 0.85) (2.30, 0.84) (2.40, 0.83)
    (2.50, 0.82)
};

\legend{Stromatolit (biyotik), Smektit kil (bağlam), Demir oksit (abiyotik), Evaporit (abiyotik)}

% Absorpsiyon bant notları
\draw[-{Stealth}, thin, gray] (axis cs:1.40, 0.38) -- (axis cs:1.40, 0.10)
  node[below, font=\tiny, text=gray!70!black] {OH/H$_2$O};
\draw[-{Stealth}, thin, gray] (axis cs:1.90, 0.25) -- (axis cs:1.90, 0.10)
  node[below, font=\tiny, text=gray!70!black] {H$_2$O};
\draw[-{Stealth}, thin, gray] (axis cs:2.30, 0.30) -- (axis cs:2.30, 0.10)
  node[below, font=\tiny, text=gray!70!black] {CO$_3$/Metal-OH};
\draw[-{Stealth}, thin, gray] (axis cs:0.90, 0.35) -- (axis cs:0.90, 0.10)
  node[below, font=\tiny, text=gray!70!black] {Fe$^{3+}$};

\end{axis}
\end{tikzpicture}
\caption{Biyotik ve abiyotik yüzey malzemelerinin tipik VNIR-SWIR yansıtma spektrumları (şematik). Stromatolit (yeşil, düz): organik madde ve karbonat absorpsiyon bantları. Smektit kil (turuncu, kesikli): 1{,}4 ve 1{,}9\,$\mu$m'de derin OH/H$_2$O bantları \cite{clark1984}. Demir oksit (kırmızı, noktalı): 0{,}9\,$\mu$m'de Fe$^{3+}$ absorpsiyonu -- biyotik yapıyı taklit edebilir. Evaporit (mor, nokta-çizgi): yüksek albedo, düz profil. Çoklu kanal yaklaşımı bu imzaları ayırt etmeyi amaçlar.}
\label{fig:spectra}
\end{figure}

% ===== ŞEKİL 9: Jezero Krateri Delta Kesiti =====
\begin{figure}[ht]
\centering
\begin{tikzpicture}[
    annot/.style={font=\tiny, text=gray!70!black},
    mlabel/.style={font=\tiny\bfseries, rounded corners=1pt, inner sep=2pt},
]
% Ölçek: yatay ~3 km, dikey ~200 m (abartılı)

% Krater tabanı
\fill[cLow!8] (0,-0.2) rectangle (14,-2.4);

% Bazalt taban katmanı
\fill[cLow!18] (0,-1.6) rectangle (14,-2.4);
\draw[cLow!40, thick] (0,-1.6) -- (14,-1.6);
\node[mlabel, fill=cLow!30] at (8.0, -1.1) {olivin bazalt};

% Ana delta gövdesi
\fill[cContext!25] (1.5, 2.5) -- (3.0, 2.8) -- (5.0, 3.0) -- (7.0, 2.8)
                   -- (9.5, 2.0) -- (11.5, 0.5) -- (12.5, -0.2)
                   -- (0.5, -0.2) -- (0.5, 1.5) -- cycle;

% Kil katmanı (smektit — alt)
\fill[cChem!20] (1.0, 0.0) -- (12.0, -0.2) -- (12.0, 0.5) -- (10.5, 1.0)
                -- (8.0, 1.5) -- (5.0, 1.8) -- (2.0, 1.5) -- (1.0, 0.8) -- cycle;
\draw[cChem!50, thick, dashed] (1.0, 0.8) -- (2.0, 1.5) -- (5.0, 1.8) -- (8.0, 1.5) -- (10.5, 1.0) -- (12.0, 0.5);

% Karbonat katmanı (üst)
\fill[cContext!40] (1.5, 2.5) -- (3.0, 2.8) -- (5.0, 3.0) -- (7.0, 2.8)
                   -- (9.0, 2.2) -- (8.0, 1.5) -- (5.0, 1.8) -- (2.0, 1.5) -- (1.0, 0.8) -- (0.8, 1.5) -- cycle;

% Delta üst yüzey çizgisi
\draw[cContext!70, very thick] (0.5, 1.5) -- (1.5, 2.5) -- (3.0, 2.8) -- (5.0, 3.0) -- (7.0, 2.8) -- (9.5, 2.0) -- (11.5, 0.5) -- (12.5, -0.2);

% Foreset katmanlar (delta öncülleri)
\foreach \x in {4.0, 5.5, 7.0, 8.5, 10.0} {
  \draw[cContext!45, thin] (\x, {2.8 - (\x-4)*0.28}) -- ({\x+1.0}, {1.0 - (\x-4)*0.15});
}

% Göl seviyesi (eski)
\draw[cRemote, thick, densely dashed] (0, 2.6) -- (14, 2.6);
\node[annot, anchor=east, text=cRemote] at (-0.15, 2.6) {\textit{paleo-göl sev.}};

% Krater kenarı (sol)
\fill[cLow!20] (0, 4.2) -- (0.3, 3.5) -- (0.5, 1.5) -- (0, 1.5) -- cycle;
\draw[cLow!60, very thick] (0, 4.2) -- (0.3, 3.5) -- (0.5, 1.5);
\node[annot, fill=white, inner sep=2pt, rounded corners=1pt, anchor=east, font=\tiny\bfseries] at (-0.15, 3.2) {krater kenarı};

% Kanal (giriş)
\fill[cRemote!30] (0, 3.5) -- (0.5, 2.8) -- (1.5, 2.5) -- (1.0, 2.7) -- (0, 3.3) -- cycle;
\draw[cRemote!60, thick] (0, 3.5) -- (0.5, 2.8) -- (1.5, 2.5);
\draw[cRemote!60, thick] (0, 3.3) -- (1.0, 2.7);
\node[annot, anchor=south west, fill=white, inner sep=2pt, rounded corners=1pt, font=\tiny\bfseries] at (0.8, 2.85) {giriş kanalı};

% Krater kenarı (sağ)
\fill[cLow!20] (14, 3.0) -- (13.5, 1.0) -- (12.5, -0.2) -- (14, -0.2) -- cycle;
\draw[cLow!60, very thick] (14, 3.0) -- (13.5, 1.0) -- (12.5, -0.2);

% Mineral etiketleri
\node[mlabel, fill=cContext!50, text=white] at (4.0, 2.4) {karbonat};
\node[mlabel, fill=cChem!50, text=white] at (6.0, 0.8) {smektit (kil)};

% Perseverance konumu
\node[star, star points=5, star point height=0.15cm, fill=cFusion, draw=cFusion!80!black, minimum size=0.3cm] (rover) at (5.0, 3.15) {};
\node[annot, anchor=south, font=\tiny\bfseries, text=cFusion!80!black] at (rover.north) {Perseverance};

% Silika zenginleşme zonu
\fill[white, opacity=0.4] (9.0, 0.3) ellipse (0.8 and 0.4);
\draw[cExtra, thick, dashed] (9.0, 0.3) ellipse (0.8 and 0.4);
\node[annot, text=cExtra] at (9.0, 0.3) {SiO$_2$};

% Ölçek çubuğu
\draw[|-|, very thick] (10.5, -2.7) -- (13.5, -2.7);
\node[font=\scriptsize, below] at (12.0, -2.7) {$\sim$1 km};

% Dikey ölçek
\draw[|-|, very thick] (14.3, -0.2) -- (14.3, 3.0);
\node[font=\scriptsize, right] at (14.3, 1.4) {$\sim$200\,m};

\end{tikzpicture}
\caption{Jezero krateri delta yapısının şematik kesiti (dikey abartılı). Giriş kanalından taşınan sediman, paleolaküstrin ortamda karbonat (üst, turuncu) ve smektit kil (alt, turkuaz) katmanları olarak çökelmiştir. Silika (SiO$_2$) zenginleşme zonları hidrotermal alterasyona işaret eder. Perseverance gezgini (yıldız) delta üst yüzeyinde örnekleme yapmaktadır \cite{farley2020, bhartia2021, allwood2020}. ``Foreset'' eğimli katmanlar delta ilerlemesinin jeolojik kaydını oluşturur.}
\label{fig:jezero}
\end{figure}

\section{Öznitelik Mühendisliği ve Model Mimarisi}
\subsection{Sensörden Türetilen Öznitelikler}
\begin{table}[ht]
\centering
\caption{Sensör kanallarına göre türetilen öznitelik örnekleri.}
\label{tab:features}
\begin{tabular}{>{\raggedright\arraybackslash}p{0.28\textwidth} >{\raggedright\arraybackslash}p{0.62\textwidth}}
\toprule
\textbf{Kanal} & \textbf{Örnek Öznitelikler} \\
\midrule
Orbital Spektral & Bant derinlik oranları, spektral eğim, absorpsiyon merkezi kayması, hidratasyon göstergeleri \\
Jeomorfoloji/Topografya & Eğim, pürüzlülük, drenaj yoğunluğu, delta/yelpaze parametreleri, katman süreklilik indeksi \\
Yüzey Görüntüleme & Fraktal boyut, yönlenme dağılımı, laminasyon periyodisitesi, filament benzerlik skoru \\
Kimyasal/Mineralojik & Organik pik yoğunluğu, fonksiyonel grup skorları, kil-sülfat-silika bağlam uyumu, redoks indeksleri \\
İzotopik & $\delta^{13}C$, $\delta^{34}S$, $\delta D$ sapmaları ve fraksiyonasyon tutarlılığı \\
Kalite/Kontaminasyon & \textit{Blank-witness} korelasyonu, numune zinciri güven puanı, cihaz kayma sapması \\
\bottomrule
\end{tabular}
\end{table}

\subsection{Önerilen Katmanlı Mimari}
\begin{enumerate}[label=\textbf{Katman \Alph*:},leftmargin=2.6cm]
    \item Aday tarama: 1D-CNN/\textit{Transformer} $\rightarrow S_{remote}$
    \item Bağlam değerlendirme: XGBoost/Bayes ağı $\rightarrow S_{context}$
    \item Morfoloji incelemesi: Çok ölçekli CNN $\rightarrow S_{in\text{-}situ}$
    \item Kimya-İzotop doğrulama: Geç birleştirme topluluğu $\rightarrow S_{chem\text{-}iso}$
    \item Karar füzyonu: Bayes güncelleme + risk ceza terimi $\rightarrow S_{final}$
\end{enumerate}

\section{Skor Kartı ve Karar Eşikleri}
\subsection{Skor Kartı Alanları}
\begin{itemize}
    \item Kimlik: bölge kimliği, koordinat, veri sürümü, tarih
    \item Ana skorlar: $S_{remote}, S_{context}, S_{in\text{-}situ}, S_{chem\text{-}iso}, R_{contam}, S_{final}$
    \item Güven: güven aralığı, belirsizlik seviyesi, karar etiketi
    \item Kanıt özeti: en güçlü pozitifler, kritik karşı kanıtlar, kontaminasyon sonuçları
\end{itemize}

\subsection{Örnek Karar Eşikleri}
\begin{center}
\begin{tabular}{lll}
\toprule
\textbf{Skor Aralığı} & \textbf{Etiket} & \textbf{Operasyon} \\
\midrule
$S_{final} < 0.35$ & Düşük güven & Eleme veya düşük öncelik \\
$0.35 \leq S_{final} < 0.60$ & Ek örnekleme gerekli & Ek ölçüm ve tekrar doğrulama \\
$0.60 \leq S_{final} < 0.80$ & Orta güven & Derin örnekleme + izotop tekrarı \\
$S_{final} \geq 0.80$ & Yüksek güven & Öncelikli bilimsel hedef \\
\bottomrule
\end{tabular}
\end{center}

% ===== ŞEKİL 2: Karar Eşikleri Görsel =====
\begin{figure}[ht]
\centering
\begin{tikzpicture}
    \def\barW{13.5}
    \def\barH{1.25}
  % Arka plan bantları
  \fill[cLow!30]   (0,0) rectangle ({0.35*\barW},\barH);
  \fill[cExtra!30]  ({0.35*\barW},0) rectangle ({0.60*\barW},\barH);
  \fill[cHigh!25]   ({0.60*\barW},0) rectangle ({0.80*\barW},\barH);
  \fill[cHigh!50]   ({0.80*\barW},0) rectangle (\barW,\barH);
  % Çerçeve
  \draw[thick] (0,0) rectangle (\barW,\barH);
  % Dikey çizgiler
    \draw[thick, dashed] ({0.35*\barW},0) -- ({0.35*\barW},{\barH+0.38});
    \draw[thick, dashed] ({0.60*\barW},0) -- ({0.60*\barW},{\barH+0.38});
    \draw[thick, dashed] ({0.80*\barW},0) -- ({0.80*\barW},{\barH+0.38});
  % Etiketler
    \node[font=\scriptsize\bfseries, text=cLow, align=center]   at ({0.175*\barW},{\barH/2}) {Düşük\\güven};
    \node[font=\scriptsize\bfseries, text=cExtra, align=center]  at ({0.475*\barW},{\barH/2}) {Ek\\örnekleme};
    \node[font=\scriptsize\bfseries, text=cHigh!80!black, align=center]  at ({0.70*\barW},{\barH/2}) {Orta\\güven};
    \node[font=\scriptsize\bfseries, text=cHigh!50!black, align=center]  at ({0.90*\barW},{\barH/2}) {Yüksek\\güven};
  % Eşik değerleri
    \node[font=\scriptsize, below] at ({0.35*\barW},-0.18)  {0.35};
    \node[font=\scriptsize, below] at ({0.60*\barW},-0.18)  {0.60};
    \node[font=\scriptsize, below] at ({0.80*\barW},-0.18)  {0.80};
    \node[font=\scriptsize, below] at (0,-0.18)  {0.00};
    \node[font=\scriptsize, below] at (\barW,-0.18) {1.00};
  % ok + baslik
    \draw[-{Stealth}, thick] (0,-0.90) -- (\barW,-0.90) node[midway, below, font=\footnotesize]{$S_{final}$};
\end{tikzpicture}
\caption{Karar eşik bölgelerinin görsel temsili.}
\label{fig:thresholds}
\end{figure}

\section{90 Günlük Asgari Ürün Yol Haritası}
\begin{enumerate}[label=\arabic*.]
    \item Gün 1--30: veri şeması, etiket standardı, ilk uzaktan tarama modeli.
    \item Gün 31--60: bağlam + morfoloji model entegrasyonu, negatif sınıf genişletme.
    \item Gün 61--90: kimyasal-izotopik füzyon, kalibrasyon, eşik optimizasyonu, gösterge paneli.
\end{enumerate}

\section{Değerlendirme ve Başarı Ölçütleri}
\begin{enumerate}[label=\arabic*.]
    \item Yanlış pozitif oranının hedef eşik altına indirilmesi,
    \item Farklı sahalarda tekrarlanabilir skor davranışı,
    \item Çoklu kanıt füzyonunun tek kanala göre anlamlı başarım artışı,
    \item Belirsizlik raporlamasıyla birlikte eyleme dönük karar üretilmesi.
\end{enumerate}

\section{Doğrulama Sonuçları}

Karar modeli, 12 referans senaryodan oluşan bir test takımı ile doğrulanmıştır. Senaryolar: Dünya pozitif kontrolleri (Shark Bay stromatoliti \cite{allwood2006}, Yellowstone sinter), Dünya negatif kontrolleri (evaporit, demir oksit, Fischer--Tropsch sentetik \cite{mccollom2006}), Mars senaryoları (Jezero delta \cite{farley2020}, sülfat damar, bazalt) ve tek kanal baskınlık testlerini içerir.

\subsection{Senaryo Sonuç Tablosu}

% ===== ŞEKİL 3: Doğrulama Çubuk Grafiği =====
\begin{figure}[ht]
\centering
\begin{tikzpicture}
\begin{axis}[
    xbar,
    width=12cm, height=12cm,
    xlabel={$S_{final}$},
    xmin=0, xmax=1.0,
    ytick=data,
    yticklabels={
        Eksik kanal cezası,
        Tek kanal (yerinde),
        Tek kanal (uzaktan),
        Mars bazalt,
        Mars sülfat damar,
        Jezero delta,
        Fischer-Tropsch,
        Demir oksit,
        Death Valley evaporit,
        Arsenik mat,
        Yellowstone sinter,
        Shark Bay stromatolit
    },
    y dir=reverse,
    ytick style={draw=none},
    yticklabel style={font=\scriptsize},
    xticklabel style={font=\scriptsize},
    bar width=10pt,
    nodes near coords,
    nodes near coords style={font=\tiny, anchor=west},
    every node near coord/.append style={xshift=2pt},
    grid=major,
    grid style={gray!20},
    legend style={at={(0.97,0.03)}, anchor=south east, font=\tiny},
    % Eşik referans çizgileri
    extra x ticks={0.35, 0.60, 0.80},
    extra x tick labels={},
    extra x tick style={grid=major, grid style={dashed, black!40}},
    enlarge y limits={abs=0.6},
]
% Gerçek validation_suite.py --json çıktısı (2026-03-25)
\addplot[fill=cHigh!70, draw=cHigh!90!black, bar shift=0pt] coordinates {
    (0.8990,0) (0.8002,1)
};
\addplot[fill=cContext!70, draw=cContext!90!black, bar shift=0pt] coordinates {
    (0.6760,2)
};
\addplot[fill=cLow!60, draw=cLow!90!black, bar shift=0pt] coordinates {
    (0.0000,3) (0.0000,4)
};
\addplot[fill=cExtra!60, draw=cExtra!90!black, bar shift=0pt] coordinates {
    (0.2582,5) (0.4417,7) (0.3247,9) (0.3018,10)
};
\addplot[fill=cContext!50, draw=cContext!80!black, bar shift=0pt] coordinates {
    (0.5767,6) (0.5317,11)
};
\addplot[fill=cLow!60, draw=cLow!90!black, bar shift=0pt] coordinates {
    (0.0222,8)
};
\legend{Yüksek güven, Orta güven, Düşük güven, Ek örnekleme, Orta/Ek sınır}
\end{axis}
\end{tikzpicture}
\caption{12 doğrulama senaryosunun $S_{final}$ skorları (gerçek \texttt{validation\_suite.py} çıktısı). Kesikli çizgiler karar eşiklerini gösterir (0{,}35, 0{,}60, 0{,}80).}
\label{fig:validation}
\end{figure}

\begin{center}
\scriptsize
\begin{tabular}{llllr}
\toprule
\textbf{Senaryo} & \textbf{Beklenen} & \textbf{Gerçek} & \textbf{MC\footnote{Monte Carlo} Kararl.} & $S_{final}$ \\
\midrule
Shark Bay stromatoliti           & Yüksek güven           & Yüksek güven           & 1,000 & 0,8990 \\
Yellowstone sinter               & Yüksek güven           & Yüksek güven           & 0,996 & 0,8002 \\
Arsenik mikrobiyal mat           & Orta güven             & Orta güven             & 0,962 & 0,6760 \\
Death Valley evaporit            & Düşük güven            & Düşük güven            & 1,000 & 0,0000 \\
Demir oksit taklitçi             & Düşük güven            & Düşük güven            & 0,984 & 0,0000 \\
Fischer--Tropsch sentetik        & Ek örnekleme gerekli   & Ek örnekleme gerekli   & 0,888 & 0,2582 \\
Jezero delta (güçlü aday)        & Orta güven             & Orta güven             & 1,000 & 0,5767 \\
Mars sülfat damar                & Ek örnekleme gerekli   & Ek örnekleme gerekli   & 0,596 & 0,4417 \\
Mars bazalt (sinyal yok)         & Düşük güven            & Düşük güven            & 1,000 & 0,0222 \\
Tek kanal -- yalnız S\_remote      & Ek örnekleme gerekli   & Ek örnekleme gerekli   & 1,000 & 0,3247 \\
Tek kanal -- yalnız S\_in\_situ    & Ek örnekleme gerekli   & Ek örnekleme gerekli   & 1,000 & 0,3018 \\
Eksik kanal cezası (3 eksik)     & Orta güven             & Orta güven             & 0,996 & 0,5317 \\
\midrule
\multicolumn{5}{l}{\textbf{Başarı: 12/12 (\%100) -- Kaynak: \texttt{validation\_suite.py --json} (2026-03-25)}} \\
\bottomrule
\end{tabular}
\end{center}

% ===== ŞEKİL 6: Karışıklık Matrisi =====
\begin{figure}[ht]
\centering
\begin{tikzpicture}
% 4x4 karışıklık matrisi -- gerçek validation_suite.py sonuçları (12/12)
% Sınıflar: Yüksek güven (2), Orta güven (3), Ek örnekleme (4), Düşük güven (3)
% Köşegen: 2, 3, 4, 3  Köşegen dışı: hepsi 0
% Elle hücre çizimi (babel \ifnum uyumsuzluğu nedeniyle)

% Satır 0: Yüksek güven => [2, 0, 0, 0]
\fill[cHigh!50!white] (0, 0) rectangle (1,-1);
\fill[gray!5]         (1, 0) rectangle (2,-1);
\fill[gray!5]         (2, 0) rectangle (3,-1);
\fill[gray!5]         (3, 0) rectangle (4,-1);
% Satır 1: Orta güven => [0, 3, 0, 0]
\fill[gray!5]         (0,-1) rectangle (1,-2);
\fill[cHigh!75!white] (1,-1) rectangle (2,-2);
\fill[gray!5]         (2,-1) rectangle (3,-2);
\fill[gray!5]         (3,-1) rectangle (4,-2);
% Satır 2: Ek örnekleme => [0, 0, 4, 0]
\fill[gray!5]         (0,-2) rectangle (1,-3);
\fill[gray!5]         (1,-2) rectangle (2,-3);
\fill[cHigh]          (2,-2) rectangle (3,-3);
\fill[gray!5]         (3,-2) rectangle (4,-3);
% Satır 3: Düşük güven => [0, 0, 0, 3]
\fill[gray!5]         (0,-3) rectangle (1,-4);
\fill[gray!5]         (1,-3) rectangle (2,-4);
\fill[gray!5]         (2,-3) rectangle (3,-4);
\fill[cHigh!75!white] (3,-3) rectangle (4,-4);

% Izgara çizgileri
\draw[gray!40] (0,0) grid (4,-4);

% Değerler
\node[font=\large\bfseries] at (0.5,-0.5) {2};
\node[font=\small, gray]    at (1.5,-0.5) {0};
\node[font=\small, gray]    at (2.5,-0.5) {0};
\node[font=\small, gray]    at (3.5,-0.5) {0};
\node[font=\small, gray]    at (0.5,-1.5) {0};
\node[font=\large\bfseries] at (1.5,-1.5) {3};
\node[font=\small, gray]    at (2.5,-1.5) {0};
\node[font=\small, gray]    at (3.5,-1.5) {0};
\node[font=\small, gray]    at (0.5,-2.5) {0};
\node[font=\small, gray]    at (1.5,-2.5) {0};
\node[font=\large\bfseries] at (2.5,-2.5) {4};
\node[font=\small, gray]    at (3.5,-2.5) {0};
\node[font=\small, gray]    at (0.5,-3.5) {0};
\node[font=\small, gray]    at (1.5,-3.5) {0};
\node[font=\small, gray]    at (2.5,-3.5) {0};
\node[font=\large\bfseries] at (3.5,-3.5) {3};

% Ust eksen: Tahmin edilen etiketler
\node[font=\scriptsize, rotate=45, anchor=south west] at (0.5, 0.1) {Yüksek};
\node[font=\scriptsize, rotate=45, anchor=south west] at (1.5, 0.1) {Orta};
\node[font=\scriptsize, rotate=45, anchor=south west] at (2.5, 0.1) {Ek örnekl.};
\node[font=\scriptsize, rotate=45, anchor=south west] at (3.5, 0.1) {Düşük};

% Sol eksen: Beklenen etiketler
\node[font=\scriptsize, anchor=east] at (-0.1,-0.5) {Yüksek};
\node[font=\scriptsize, anchor=east] at (-0.1,-1.5) {Orta};
\node[font=\scriptsize, anchor=east] at (-0.1,-2.5) {Ek örnekl.};
\node[font=\scriptsize, anchor=east] at (-0.1,-3.5) {Düşük};

% Başlıklar
\node[font=\small, anchor=south] at (2, 1.6) {Tahmin edilen etiket};
\node[font=\small, anchor=south, rotate=90] at (-1.8, -2) {Beklenen etiket};
\end{tikzpicture}
\caption{Karışıklık matrisi -- 12 doğrulama senaryosu. Tam köşegen eşleşme (\%100 doğruluk). Gerçek \texttt{validation\_suite.py} çıktısı.}
\label{fig:confusion}
\end{figure}

\subsection{Monte Carlo Belirsizlik Analizi}

Orta güven senaryosu ($S_{final} = 0{,}7589$, tüm analoglar güçlü) üzerinde $N = 2000$ Monte Carlo pertürbasyonu ($\sigma = 0{,}05$, tohum=42) uygulanmıştır \cite{metropolis1949}:

\begin{itemize}
    \item Ortalama: 0{,}7589, Standart sapma: 0{,}0297
    \item \%90 güven aralığı: $[0{,}7096,\; 0{,}8078]$
    \item Etiket dağılımı: Yüksek güven \%87{,}25, Orta güven \%10{,}90, Ek örnekleme \%1{,}85
\end{itemize}

Bu sonuç, orta güven--yüksek güven karar sınırının (0{,}80) yakınında etiket geçişi olduğunu göstermektedir.

Ek olarak, gerçek işlem hattı çıktısı üzerinde ($S_{final} = 0{,}3976$, spektral+izotop verilerinden) $N = 1000$ MC testi yapılmıştır:
\begin{itemize}
    \item Ortalama: 0{,}4002, Standart sapma: 0{,}0246
    \item \%90 güven aralığı: $[0{,}3603,\; 0{,}4416]$
    \item Etiket dağılımı: Ek örnekleme gerekli \%97{,}6, Orta güven \%2{,}4
\end{itemize}

% ===== ŞEKİL 4: Monte Carlo Histogram =====
\begin{figure}[ht]
\centering
\begin{tikzpicture}
\begin{axis}[
    ybar,
    width=10cm, height=6cm,
    xlabel={$S_{final}$},
    ylabel={Frekans},
    xmin=0.64, xmax=0.88,
    ymin=0, ymax=420,
    xtick={0.66,0.70,0.74,0.78,0.82,0.86},
    bar width=6pt,
    grid=major,
    grid style={gray!15},
    xticklabel style={font=\scriptsize},
    yticklabel style={font=\scriptsize},
    xlabel style={font=\small},
    ylabel style={font=\small},
    fill=cRemote!60,
    draw=cRemote!80!black,
]
% Gerçek MC dağılımı: score_evidence.py N=2000, sigma=0.05, tohum=42
% Referans senaryo: ortalama=0.7589, std=0.0297
\addplot coordinates {
    (0.6600,   1) (0.6750,  13) (0.6900,  52) (0.7050, 152)
    (0.7200, 221) (0.7350, 335) (0.7500, 365) (0.7650, 369)
    (0.7800, 269) (0.7950, 134) (0.8100,  68) (0.8250,  18)
    (0.8400,   2) (0.8550,   1)
};
% GA90 çizgileri (gerçek P5=0.7096, P95=0.8078)
\draw[thick, dashed, cLow] (axis cs:0.7096,0) -- (axis cs:0.7096,390) node[above, font=\tiny]{P5=0.710};
\draw[thick, dashed, cLow] (axis cs:0.8078,0) -- (axis cs:0.8078,390) node[above, font=\tiny]{P95=0.808};
% Ortalama
\draw[thick, cHigh] (axis cs:0.7589,0) -- (axis cs:0.7589,390) node[above, font=\tiny]{$\mu$=0.759};
\end{axis}
\end{tikzpicture}
\caption{Monte Carlo belirsizlik dağılımı ($N=2000$, $\sigma=0{,}05$, tohum=42). Gerçek \texttt{score\_evidence.py} çıktısı. Kırmızı kesikli çizgiler \%90 güven aralığını ($[0{,}710,\; 0{,}808]$), yeşil çizgi ortalama skoru ($\mu=0{,}759$) gösterir.}
\label{fig:montecarlo}
\end{figure}

\subsection{Ağırlık Hassasiyet Analizi}

Her ağırlık $\pm$\%20 pertürbasyona tabi tutulmuştur (Tablo~\ref{tab:sensitivity}). Referans senaryoda en hassas parametre $\delta$ (S\_chem\_iso katsayısı) olup $\pm 0{,}063$ skor değişimi üretmiştir. En az hassas parametre $\lambda$ (kontaminasyon katsayısı, $\pm 0{,}007$) bulunmuştur.

\begin{table}[ht]
\centering
\caption{Ağırlık hassasiyet analizi -- referans senaryo (gerçek \texttt{score\_evidence.py} çıktısı).}
\label{tab:sensitivity}
\begin{tabular}{lcrrr}
\toprule
\textbf{Ağırlık} & \textbf{Varsayılan} & \textbf{$-$20\%} & \textbf{$+$20\%} & $|\Delta S|_{max}$ \\
\midrule
$\alpha$ (S\_remote)   & 0.25 & $-$0.0387 & $+$0.0386 & 0.0387 \\
$\beta$ (S\_context)   & 0.20 & $-$0.0258 & $+$0.0257 & 0.0258 \\
$\gamma$ (S\_in\_situ) & 0.20 & $-$0.0312 & $+$0.0311 & 0.0312 \\
$\delta$ (S\_chem\_iso)& 0.35 & $-$0.0630 & $+$0.0629 & 0.0630 \\
$\lambda$ (R\_contam)  & 0.25 & $+$0.0067 & $-$0.0068 & 0.0068 \\
\bottomrule
\end{tabular}
\end{table}

Bu durum, kimyasal-izotopik doğrulama kanalının ($\delta$) modelin nihai kararında belirleyici rol oynadığını doğrulamaktadır.

% ===== ŞEKİL 5: Hassasiyet Çubuk Grafiği =====
\begin{figure}[ht]
\centering
\begin{tikzpicture}
\begin{axis}[
    ybar,
    width=10cm, height=5.5cm,
    ylabel={$|\Delta S_{final}|$},
    symbolic x coords={$\alpha$, $\beta$, $\gamma$, $\delta$, $\lambda$},
    xtick=data,
    xticklabel style={font=\small},
    yticklabel style={font=\scriptsize},
    ylabel style={font=\small},
    ymin=0, ymax=0.075,
    bar width=18pt,
    nodes near coords,
    nodes near coords style={font=\scriptsize},
    grid=major,
    grid style={gray!15},
    every node near coord/.append style={anchor=south},
]
% Gerçek sensitivity_analysis() çıktısı (referans senaryo)
\addplot[fill=cRemote!60, draw=cRemote!80!black] coordinates {
    ($\alpha$, 0.0387)
    ($\beta$,  0.0258)
    ($\gamma$, 0.0312)
    ($\delta$, 0.0630)
    ($\lambda$,0.0068)
};
\end{axis}
\end{tikzpicture}
\caption{Ağırlık hassasiyet analizi ($\pm$\%20 pertürbasyon). Gerçek \texttt{score\_evidence.py} çıktısı. $\delta$ (kimyasal-izotopik kanal) en hassas parametredir.}
\label{fig:sensitivity}
\end{figure}

\section{Sonuç}
Bu makale, astrobiyoloji görevlerinde ``uzaktan tarama + yakın doğrulama'' mantığını tek bir biyo-iz güven skorunda birleştiren uygulanabilir bir çerçeve sunmaktadır. Önerilen model, bilimsel güveni artırırken operasyonel kaynak kullanımını optimize eder. Özellikle Mars gibi jeolojik olarak zengin ortamlarda, çok modlu veri füzyonu sayesinde yanlış pozitiflerin azaltılması ve hedef seçiminin iyileştirilmesi beklenmektedir.

Doğrulama testleri, modelin bilinen pozitif, negatif ve belirsiz senaryolarda tutarlı karar ürettiğini göstermektedir (12/12 başarı, Şekil~\ref{fig:confusion}). Monte Carlo belirsizlik analizi, karar sınırları yakınında etiket geçişlerinin doğal olduğunu ve operasyonel olarak ``ek örnekleme'' önerisiyle yönetilebildiğini ortaya koymaktadır. Ağırlık hassasiyet analizi, kimyasal-izotopik kanalın modelin belirleyici giriş değişkeni olduğunu doğrulamakta ve gelecek kalibrasyonların bu kanal üzerine yoğunlaşmasını önermektedir.

\subsection{Sınırlılıklar}
Mevcut çalışma bazı önemli sınırlılıklar içermektedir:
\begin{enumerate}
    \item \textbf{Tek saha kalibrasyonu:} Karar eşikleri (0{,}35, 0{,}60, 0{,}80) yalnızca Shark Bay, Yellowstone ve Jezero analog senaryolarına dayalı olarak ayarlanmıştır. Enceladus, Europa veya Titan gibi farklı jeokimyasal ortamlara transferi henüz doğrulanmamıştır.
    \item \textbf{Sentetik test verileri:} Doğrulama senaryoları literatür değerlerine dayalı olarak elle oluşturulmuştur; gerçek görev verileriyle uçtan uca test henüz yapılamamıştır.
    \item \textbf{Makine öğrenmesi eksikliği:} Mevcut model kural tabanlı bir karar ağacı kullanmaktadır. Gözetimli veya gözetimsiz sınıflandırma henüz entegre edilmemiştir.
    \item \textbf{Kontaminasyon modeli basitliği:} $R_{contam}$ parametresi sabit bir ceza katsayısı olarak tanımlanmıştır; gerçek kontaminasyon senaryoları için olasılıksal bir model gereklidir.
    \item \textbf{Zamansal evrim eksikliği:} Model, tek bir zaman dilimindeki ölçümleri birleştirir; farklı sollarda tekrarlanan ölçümlerin Bayes güncelleme ile biriktirilmesi henüz uygulanmamıştır.
    \item \textbf{Ay doğrulama eksikliği:} Özet ve yöntem bölümlerinde Ay hedef olarak tanımlanmış olmasına karşın, mevcut doğrulama senaryoları yalnızca Dünya ve Mars sahalarını kapsamaktadır. Ay'a özgü senaryolar (örn.\ PSR buz deposu, mare bazalt) gelecek çalışmalara bırakılmıştır.
\end{enumerate}

\subsection{Gelecek Çalışmalar}
Yukarıdaki sınırlılıklar ve mevcut işlem hattı üzerine aşağıdaki geliştirmeler planlanmaktadır:
\begin{enumerate}
    \item \textbf{SDSS spektral entegrasyonu:} Teleskop spektroskopi FITS dosyalarının $S_{remote}$ kanalına otomatik beslenmesi; derin uzay biyo-iz tarama modülü.
    \item \textbf{Gerçek Mars PDS verisi ile uçtan uca test:} SHERLOC Raman+floresans, PIXL XRF ve SuperCam LIBS kalibrasyon verilerinin işlem hattına entegrasyonu (veriler bu çalışmada indirilmiştir, bkz.\ Şekil~\ref{fig:pipeline}).
    \item \textbf{Temel makine öğrenmesi sınıflandırıcısı:} Doğrulama senaryolarından öğrenen bir rastgele orman (\textit{random forest}) veya gradyan artırma sınıflandırıcısı ile kural tabanlı karar ağacının karşılaştırılması.
    \item \textbf{Çoklu görev transferi:} LRO/LOLA, KPLO/ShadowCam ve Cassini/CIRS verilerinin Ay ve dış gezegen uydu senaryolarına genişletilmesi.
    \item \textbf{Etkileşimli görselleştirme:} Web tabanlı bir gösterge paneliyle senaryo parametrelerinin canlı olarak keşfedilmesi ve karar sınırlarının dinamik güncellenmesi.
\end{enumerate}

\begin{thebibliography}{99}

% --- Mars 2020 Görevi ve Aletleri ---
\bibitem{farley2020}
Farley, K.~A. \textit{et al.},
``Mars 2020 Mission Overview,''
\textit{Space Science Reviews}, 216(8), 142, 2020.
\texttt{doi:10.1007/s11214-020-00762-y}

\bibitem{bhartia2021}
Bhartia, R. \textit{et al.},
``Perseverance's Scanning Habitable Environments with Raman \& Luminescence for Organics \& Chemicals (SHERLOC) Investigation,''
\textit{Space Science Reviews}, 217(4), 58, 2021.
\texttt{doi:10.1007/s11214-021-00812-z}

\bibitem{allwood2020}
Allwood, A.~C. \textit{et al.},
``PIXL: Planetary Instrument for X-Ray Lithochemistry,''
\textit{Space Science Reviews}, 216(8), 134, 2020.
\texttt{doi:10.1007/s11214-020-00767-7}

\bibitem{maurice2021}
Maurice, S. \textit{et al.},
``The SuperCam Instrument Suite on the Mars 2020 Rover: Science Objectives and Mast-Unit Description,''
\textit{Space Science Reviews}, 217(3), 47, 2021.
\texttt{doi:10.1007/s11214-021-00807-w}

% --- MRO Orbital Aletleri ---
\bibitem{murchie2007}
Murchie, S. \textit{et al.},
``Compact Reconnaissance Imaging Spectrometer for Mars (CRISM) on Mars Reconnaissance Orbiter (MRO),''
\textit{Journal of Geophysical Research}, 112(E5), E05S03, 2007.
\texttt{doi:10.1029/2006JE002682}

\bibitem{mcewen2007}
McEwen, A.~S. \textit{et al.},
``Mars Reconnaissance Orbiter's High Resolution Imaging Science Experiment (HiRISE),''
\textit{Journal of Geophysical Research}, 112(E5), E05S02, 2007.
\texttt{doi:10.1029/2005JE002605}

\bibitem{seu2007}
Seu, R. \textit{et al.},
``SHARAD sounding radar on the Mars Reconnaissance Orbiter,''
\textit{Journal of Geophysical Research}, 112(E5), E05S05, 2007.
\texttt{doi:10.1029/2006JE002745}

% --- Ay Veri Kaynakları ---
\bibitem{smith2010}
Smith, D.~E. \textit{et al.},
``The Lunar Orbiter Laser Altimeter Investigation on the Lunar Reconnaissance Orbiter Mission,''
\textit{Space Science Reviews}, 150(1--4), 209--241, 2010.
\texttt{doi:10.1007/s11214-009-9512-y}

\bibitem{robinson2010}
Robinson, M.~S. \textit{et al.},
``Lunar Reconnaissance Orbiter Camera (LROC) Instrument Overview,''
\textit{Space Science Reviews}, 150(1--4), 81--124, 2010.
\texttt{doi:10.1007/s11214-010-9634-2}

% --- Dünya Analog ve Spektral Referanslar ---
\bibitem{kokaly2017}
Kokaly, R.~F. \textit{et al.},
``USGS Spectral Library Version 7,''
\textit{U.S. Geological Survey Data Series}, 1035, 2017.
\texttt{doi:10.5066/F7RR1WDJ}

\bibitem{sdssdr17}
Abdurro'uf \textit{et al.},
``The Seventeenth Data Release of the Sloan Digital Sky Surveys: Complete Release of MaNGA, MaStar and BOSS/eBOSS,''
\textit{The Astrophysical Journal Supplement Series}, 259(2), 35, 2022.
\texttt{doi:10.3847/1538-4365/ac4414}

% --- Spektral Analiz Yöntemleri ---
\bibitem{clark1984}
Clark, R.~N. and Roush, T.~L.,
``Reflectance Spectroscopy: Quantitative Analysis Techniques for Remote Sensing Applications,''
\textit{Journal of Geophysical Research}, 89(B7), 6329--6340, 1984.
\texttt{doi:10.1029/JB089iB07p06329}

% --- İzotop Biyo-izleri ---
\bibitem{schidlowski2001}
Schidlowski, M.,
``Carbon isotopes as biogeochemical recorders of life over 3.8~Ga of Earth history: evolution of a concept,''
\textit{Precambrian Research}, 106(1--2), 117--134, 2001.
\texttt{doi:10.1016/S0301-9268(00)00128-5}

% --- Stromatolit ve Abiyotik Kontrol ---
\bibitem{allwood2006}
Allwood, A.~C. \textit{et al.},
``Stromatolite reef from the Early Archaean era of Australia,''
\textit{Nature}, 441, 714--718, 2006.
\texttt{doi:10.1038/nature04764}

\bibitem{mccollom2006}
McCollom, T.~M. and Seewald, J.~S.,
``Carbon isotope composition of organic compounds produced by abiotic synthesis under hydrothermal conditions,''
\textit{Earth and Planetary Science Letters}, 243(1--2), 74--84, 2006.
\texttt{doi:10.1016/j.epsl.2006.01.027}

% --- NASA Veri Arşivleri ---
\bibitem{pds}
NASA Planetary Data System,
\\\texttt{https://pds.nasa.gov/}
(Erişim: 29 Mart 2026).

\bibitem{earthdata}
NASA Earthdata,
\\\texttt{https://www.earthdata.nasa.gov/}
(Erişim: 29 Mart 2026).

% --- Belirsizlik ve Monte Carlo ---
\bibitem{metropolis1949}
Metropolis, N. and Ulam, S.,
``The Monte Carlo Method,''
\textit{Journal of the American Statistical Association}, 44(247), 335--341, 1949.
\texttt{doi:10.1080/01621459.1949.10483310}

\end{thebibliography}

\subsection*{Yapay Zekâ Kullanım Bildirimi}
Bu çalışmanın hazırlanmasında Anthropic Claude Opus 4.6 dil modelinden yazım, düzenleme ve kodlama desteği alınmıştır. Tüm bilimsel içerik, yorum ve kararların nihai sorumluluğu yazarlara aittir.

\end{document}
